\section{Continual Learning Paradigms: CPT, CIT, and CA}

Continual learning in large language models is structured around a three-stage alignment framework \parencite{wu2024continuallearninglargelanguage}:
\begin{enumerate}
    \item \textbf{Continual Pre-Training (CPT):} Employs self-supervised learning on raw text corpora to inject new domain vocabulary, facts, and linguistic entities.
    \item \textbf{Continual Instruction Tuning (CIT):} Utilizes supervised fine-tuning (SFT) on prompt-response pairs to convert static parameters into active task-solving skills.
    \item \textbf{Continual Alignment (CA):} Applies reinforcement learning from human feedback (RLHF) or direct preference optimization (DPO) to align outputs with safety, policy, and ethical guidelines.
\end{enumerate}
For the Vietnamese legal domain, this alignment pipeline maps to pre-training on statutory codes, tuning on legal reasoning tasks, and preference alignment to enforce strict judicial compliance.

\subsection{Cross-Stage Forgetting and Forgetting Dynamics}
A primary vulnerability of sequential optimization is cross-stage forgetting, where optimizations in downstream stages (CIT/CA) override and corrupt knowledge acquired during early stages (CPT) \parencite{wu2024continuallearninglargelanguage}. Consequently, establishing cross-stage evaluation protocols is vital to verify that legal factual integrity is not compromised during alignment.

From a dynamical perspective, forgetting is governed by the Inverse Linear Law, establishing that task adaptation efficiency correlates with historical knowledge decay \parencite{kalajdzievski2024scalinglawsforgettingfinetuning}. Geometrically, training on out-of-distribution (OOD) data pushes model weights into sharp minima of the loss landscape, which are highly sensitive to minor parameter updates \parencite{li2024revisitingcatastrophicforgettinglarge}. Sharpness-Aware Minimization (SAM) is often utilized to locate flat minima, mitigating parameter sensitivity. Empirically, decoder-only architectures exhibit stronger resistance to catastrophic forgetting than encoder-decoder models \parencite{luo2025empiricalstudycatastrophicforgetting}.

\subsection{The Role of Replay and Rehearsal}
Data rehearsal (replay) is a standard baseline to stabilize training landscapes in CL. For CIT, replaying a minor fraction ($0.25\% - 1\%$) of historical instructions is sufficient to maintain model capabilities \parencite{scialom-etal-2022-fine}. However, for CPT, the required replay ratio scales with the severity of the domain shift: minor shifts demand roughly $5\%$ replay \parencite{ibrahim2024simplescalablestrategiescontinually}, whereas strong shifts—such as adapting to new languages or complex legal terminology—require up to $15\% - 25\%$ replay to prevent parameter drift \parencite{zheng-etal-2024-breaking}.

\subsection{CPT Optimization Mechanisms and Stability Gap}
During CPT, utilizing learning rate re-warming and re-decay schedules is mandatory to suppress loss spikes \parencite{gupta2023continualpretraininglargelanguage, yildiz2025investigatingcontinualpretraininglarge}. Additionally, transitioning across distinct data streams introduces a stability gap \parencite{du2024unlockingcontinuallearningabilities}, characterized by a temporary performance drop. This gap is primarily driven by optimizer state adjustments (momentum lag in AdamW) rather than physical data mismatch. Implementing Infinite learning rate schedules and careful warmup mitigation helps secure optimization continuity.

\subsection{Domain-Specific Scaling Laws and Evaluation Benchmarks}
Domain-specific scaling laws, such as D-CPT Law \parencite{que2024dcptlawdomainspecificcontinual}, model pre-training loss as a function of the data mixing ratio $r$. This allows researchers to project domain performance with minimal computational overhead \parencite{li2025ticlmwebscalebenchmarktimecontinual, peng2024scalablelanguagemodelgeneralized}. 

To evaluate forgetting quantitatively, benchmarks such as TRACE \parencite{wang2023tracecomprehensivebenchmarkcontinual} and ConTinTin \parencite{yin-etal-2022-contintin} measure metrics including Overall Performance (OP), Backward Transfer (BWT), Forward Transfer (FWT), and safety degradation ($\Delta R$). Evaluations indicate that while safety alignment is relatively stable, multi-step logical reasoning is highly vulnerable to decay during continual training.
